\section{Step 84 Admit the destination geometry}
Packing reads reconstructed RGB samples and writes native display words. The descriptor supplies dimensions, destination stride and writable capacity. These fields describe different quantities: width counts active pixels; stride counts words between row starts; capacity counts available bytes.

Fill \texttt{pack\_admit}. Load the descriptor’s four unsigned doublewords into \texttt{\$t5} through \texttt{\$t8}: width, height, stride and capacity. Require dimensions from one through 512 and stride from width through 512. Require a word-aligned framebuffer address. Reject through \texttt{pack\_early} before any source load or destination store. Retain the mode arriving in \texttt{\$t9} for later setup.

The worked direct-call example uses width two, height five, stride seventeen and capacity 280. The supplied file caller uses stride 512 and a 1 MiB framebuffer. Both geometries are valid; neither implies that every allocated word is part of the image. At this checkpoint, inspect the four descriptor values before expecting any pixels to change.

\paragraph{Observe} Use \texttt{s84-1}; stop before \texttt{pack\_admitted}. Inspect the state described above before continuing.

\Needspace{7\baselineskip}
\paragraph{Check} Why may stride exceed width?
\begin{itemize}
\item[A.] It describes the distance between row starts, including padding.
\item[B.] It includes compressed input bytes.
\item[C.] It changes the RGB channel order.
\end{itemize}
